(1) 绳的拉力始终沿径向，因此原点处的力矩为零，角动量 \(\ell=mr^2\dot\phi=mr v_\phi\) 守恒。径向牛顿方程 \(m(\ddot r-r\dot\phi^2)=-T\)，其中 \(T\) 为绳张力，这同时确定缓慢收绳时的准静态径向平衡。
(2) 在缓慢收绳且径向动能可忽略时，\(v_\phi(r)=\ell/(mr)\)，故 \(v_{\phi2}/v_{\phi1}=r_1/r_2\)。对应的切向动能为 \(K=\ell^2/(2mr^2)\)，增加量 \(\Delta K=\ell^2(1/r_2^2-1/r_1^2)/(2m)\)，严格为正。
(3) 准静态径向方程要求 \(T=\ell^2/(mr^3)\)，因径向位移向内，绳拉力做功 \(dW=-T\,dr\)。积分得 \(W=\int_{r_2}^{r_1}\ell^2/(mr^3)\,dr=\ell^2[1/r_2^2-1/r_1^2]/(2m)\)，与动能增量完全一致。
(4) 角动量守恒的条件是力矩为零，而不是外力做功为零；牵引力虽沿径向，却沿实际收绳位移做正功。形式极限 \(r_2\to0\) 导致速度和动能发散，实际会受孔径、绳强度与速度尺度限制，故不应外推为可行操作。
